\documentclass[11pt]{amsart}

\usepackage[utf8]{inputenc}
\usepackage[T1]{fontenc}
\usepackage{amsmath,amssymb,amsthm}
\usepackage{mathrsfs}
\usepackage{hyperref}
\usepackage{url}
\usepackage{booktabs}
\usepackage{array}
\usepackage{listings}
\usepackage{xcolor}
\usepackage[margin=1in]{geometry}

% Lean code style
\lstdefinelanguage{Lean}{
  keywords={def, theorem, instance, structure, where, by, letI, import, open, noncomputable, section, end, set_option, infer_instance, inferInstance, sorry, axiom, class, extends},
  keywordstyle=\bfseries\color{blue!70!black},
  comment=[l]{--},
  commentstyle=\itshape\color{gray},
  stringstyle=\color{red!60!black},
  basicstyle=\ttfamily\small,
  breaklines=true,
  frame=single,
  framerule=0.4pt,
  rulecolor=\color{gray!40},
  backgroundcolor=\color{gray!5},
}
\lstset{language=Lean}

% Theorem environments
\theoremstyle{plain}
\newtheorem{theorem}{Theorem}[section]
\newtheorem{proposition}[theorem]{Proposition}
\newtheorem{corollary}[theorem]{Corollary}
\newtheorem{lemma}[theorem]{Lemma}
\theoremstyle{definition}
\newtheorem{definition}[theorem]{Definition}
\newtheorem{remark}[theorem]{Remark}

% Notation
\newcommand{\Cob}{\mathrm{Cob}}
\newcommand{\Vect}{\mathrm{Vect}}
\newcommand{\TopAb}{\mathrm{TopAb}}
\newcommand{\TopVect}{\mathrm{TopVect}}
\newcommand{\LiqVect}{\mathrm{LiqVect}}
\newcommand{\CondAb}{\mathrm{CondensedAb}}
\newcommand{\CompHaus}{\mathrm{CompHaus}}
\newcommand{\Ab}{\mathrm{Ab}}
\newcommand{\Ban}{\mathrm{Ban}}
\newcommand{\SemiNormedGrp}{\mathrm{SemiNormedGrp}}
\newcommand{\ModuleCat}{\mathrm{ModuleCat}}
\newcommand{\Sheaf}{\mathrm{Sheaf}}
\newcommand{\Hom}{\mathrm{Hom}}
\newcommand{\id}{\mathrm{id}}
\newcommand{\op}{\mathrm{op}}

\title[Liquid TQFT]{Condensed Mathematics Meets Topological Quantum Field Theory:\\A Lean-Checked Exploratory Foundation}

\author{Sean Mahdavian}
\address{Phillips Exeter Academy}
\email{smhdavian@exeter.edu}

\thanks{Formal verification assisted by Aristotle (Harmonic). Mathematical synthesis and prompt engineering by Claude (Anthropic).}

\date{\today}

\begin{document}

\begin{abstract}
We present a Lean-checked investigation of condensed abelian groups as possible targets for monoidal field-theory constructions. The repository instantiates Mathlib's existing symmetric monoidal structure on $\CondAb$, defines braided monoidal composition scaffolding, constructs a faithful functor $\SemiNormedGrp \to \CondAb$, and proves that this same sheaf-level functor is not full using an explicit continuous unbounded summation map. It also formalizes commutative Frobenius algebra data and an ordinary presentation inspired by 2-dimensional cobordisms.

The project contains six active Lean files and five remaining \texttt{sorry} placeholders across three incomplete declarations, with no custom axiom declarations. It does not yet formalize a geometric bordism category, a standard Atiyah--Segal TQFT, a monoidal analytic embedding, or an LTE theorem.
\end{abstract}

\subjclass[2020]{18M05, 57R56, 18F20, 46B99, 68V15}

\keywords{condensed mathematics, topological quantum field theory, formal verification, Lean 4, monoidal categories, Frobenius algebras}

\maketitle

\tableofcontents

%% ============================================================
\section{Introduction}
\label{sec:intro}
%% ============================================================

\subsection{The Categorical Obstruction in Infinite-Dimensional TQFT}

A topological quantum field theory, in Atiyah's formulation~\cite{Atiyah1988}, is a symmetric monoidal functor $Z : \Cob_{d+1} \to C$ from a cobordism category to some symmetric monoidal target category~$C$. For finite-dimensional state spaces (Dijkgraaf--Witten, Reshetikhin--Turaev~\cite{ReshetikhinTuraev1991}, and Turaev--Viro theories), the target $C = \Vect$ works perfectly.

The physically interesting theories force infinite-dimensional state spaces. Chern--Simons theory~\cite{Witten1988} with non-compact gauge group $\mathrm{SL}(2,\mathbb{C})$, quantum gravity partition functions, and path integrals over infinite-dimensional moduli all produce state spaces $Z(\Sigma)$ that cannot live in finite-dimensional $\Vect$. The natural target becomes topological vector spaces.

Here a precise categorical obstruction arises. The category $\TopAb$ of topological abelian groups is \emph{not abelian}: a continuous bijective homomorphism need not have a continuous inverse, so a morphism can be simultaneously monic and epic without being an isomorphism. This propagates through the homological algebra machinery. Exact sequences break, derived functors misbehave, and naive homological or analytic constructions may fail. Ordinary TQFT gluing, however, is functorial composition and does not itself require exact sequences.

The completed tensor product compounds the problem. It is not exact: tensoring a short exact sequence with a topological vector space can destroy exactness. This failure can obstruct particular homological or analytic constructions involving completed tensor products. It does not, by itself, violate the ordinary symmetric-monoidal TQFT axioms.

\subsection{The Condensed Fix}

Clausen and Scholze's condensed mathematics~\cite{Scholze2019} resolves this obstruction. The category of condensed abelian groups (sheaves of abelian groups on the coherent site of compact Hausdorff spaces) is an abelian category with all limits and colimits. Every morphism has a kernel and cokernel. Every monic epic is an isomorphism.

Liquid mathematics includes special exactness and Ext-vanishing results. The Liquid Tensor Experiment~\cite{Scholze2021,LTE2022} formalized a precise Ext-vanishing theorem rather than an unrestricted exactness statement for every tensor product in $\CondAb$. No LTE theorem is reproduced here.

\subsection{This Paper}

We report on a formal verification effort investigating whether $\CondAb$ can serve as a TQFT target category. The work was conducted iteratively over six sessions using Harmonic's Aristotle system~\cite{Aristotle2025} for Lean~4 proof search, with Claude (Anthropic) providing mathematical synthesis, prompt engineering, and analysis.

The project produced both expected confirmations and genuine mathematical surprises. The monoidal structure on $\CondAb$ turned out to already exist in Mathlib (we just had to find it). The embedding from seminormed groups turned out not to be full: the obstruction was first machine-checked at the presheaf level and is now also proved for the genuine sheaf-level functor into $\CondAb$. Both findings reshaped the project's direction in real time.

%% ============================================================
\section{The Formalization: Six Sessions}
\label{sec:sessions}
%% ============================================================

\subsection{Session 1: Abstract TQFT Framework}

The first session established the categorical scaffolding in the file \texttt{LiquidTQFT.lean}.

\begin{definition}
An \emph{abstract braided monoidal theory} with source category~$S$ and target category~$C$ is a braided monoidal functor $Z : S \to C$. A standard TQFT additionally requires an actual symmetric monoidal bordism source.
\end{definition}

\begin{lstlisting}
structure AbstractTQFT
    (S : Type*) [Category S] [MonoidalCategory S] [BraidedCategory S]
    (C : Type*) [Category C] [MonoidalCategory C] [BraidedCategory C] where
  Z : S ⥤ C
  monoidal : Z.Monoidal
  braided : Z.Braided
\end{lstlisting}

We work at the level of braided (not symmetric) monoidal functors. This is reusable braided monoidal scaffolding rather than Atiyah's symmetric TQFT definition; an actual bordism source and symmetric compatibility remain to be formalized (e.g., Reshetikhin--Turaev invariants~\cite{ReshetikhinTuraev1991} from modular tensor categories). Lurie's cobordism hypothesis~\cite{Lurie2009} provides the fully extended framework in which these structures are most naturally understood.

\begin{theorem}[Transfer, sorry-free]
\label{thm:transfer}
If $T$ is an abstract braided monoidal theory valued in~$C$ and $F : C \to D$ is braided monoidal, then $T \circ F$ has the same abstract structure.
\end{theorem}

\begin{lstlisting}
def AbstractTQFT.transfer
    (T : AbstractTQFT S C)
    (F : C ⥤ D) [hF : F.Monoidal] [hFb : F.Braided] :
    AbstractTQFT S D where
  Z := T.Z ⋙ F
  monoidal := by letI := T.monoidal; infer_instance
  braided := by letI := T.monoidal; letI := T.braided; infer_instance
\end{lstlisting}

The proof is Lean's typeclass inference composing the monoidal and braided structures. Additional sorry-free results from this session: $\CondAb$ is abelian, has kernels, cokernels, and finite (co)limits, and satisfies the balanced category property (mono + epi = iso).

The session also recorded structural observations about gluing: in any abelian category, short exact sequences provide mono, epi, and exactness data. This is additional categorical structure, not ordinary TQFT gluing. These take exactness as a hypothesis. The substantive analytic claim (that the liquid tensor product is exact) depends on the LTE and is not verified in this codebase.

\subsection{Session 2: Manual Monoidal Construction}

The second session (\texttt{CondensedMonoidal.lean}) attempted to build \texttt{MonoidalCategory CondensedAb} directly, defining the tensor product as sheafification of the pointwise tensor:
\[
F \otimes_{\mathrm{cond}} G = \mathrm{sheafify}(F \otimes_{\mathrm{presheaf}} G).
\]
This yielded a complete \texttt{MonoidalCategoryStruct} (sorry-free) and 5/10 \texttt{MonoidalCategory} axioms. The remaining 5 reduced to two ``sheafification comparison'' isomorphisms:
\[
\mathrm{sheafify}(P \otimes Q) \cong \mathrm{sheafify}(\mathrm{sheafify}(P)^{\mathrm{val}} \otimes Q).
\]
This analysis was essential for the breakthrough in Session~3. As this manual construction was entirely superseded, it is not included in the final repository; we describe it here because the failed attempt is what motivated the search that found the localization approach.

\subsection{Session 3: The Localization Breakthrough}

The third session discovered that \texttt{Mathlib.CategoryTheory.Sites.Monoidal}~\cite{MatlibSitesMon} already contained the complete infrastructure.

\begin{theorem}[Sorry-free]
\label{thm:monoidal}
$\CondAb$ is a symmetric monoidal category.
\end{theorem}

The construction assembles four existing Mathlib components:
\begin{enumerate}
\item \textbf{Presheaf monoidal structure}: pointwise tensor on $\CompHaus^{\op} \to \Ab$.
\item \textbf{Sheafification as localization}: \texttt{presheafToSheaf} is a localization functor for the class~$J.W$ of local isomorphisms.
\item \textbf{$W$ is monoidal}: $J.W$ is closed under tensoring, proved via the internal hom (monoidal closed structure), not stalks.
\item \textbf{Localized monoidal category}: general construction from~\cite{MatlibLocMon}.
\end{enumerate}

\begin{lstlisting}
instance condensedAb_W_isMonoidal :
    ((coherentTopology CompHaus).W
      (A := ModuleCat (ULift ℤ))).IsMonoidal :=
  GrothendieckTopology.W.monoidal

instance condensedAb_monoidalCategory :
    MonoidalCategory CondensedAb :=
  Sheaf.monoidalCategory _ _
\end{lstlisting}

Six lines, zero sorries. The internal hom argument (rather than stalks) powers the \texttt{W.IsMonoidal} proof, meaning the result holds for any closed braided target category. A direct stalk-based approach has since been formalized by Riou (Mathlib PR \#35584), providing an alternative proof path.

\subsection{Session 4: The Banach Embedding}

The fourth session (\texttt{BanachEmbedding.lean}) did not find an existing packaged bridge between the relevant normed-group and condensed components in the pinned Mathlib version; no global priority claim is made.

\begin{definition}
For a seminormed abelian group~$V$, the presheaf $\mathrm{banachPresheaf}\;V$ sends $S \in \CompHaus$ to $C(S,V)$, the abelian group of continuous maps, with precomposition as functorial action.
\end{definition}

\begin{theorem}[Sorry-free]
\label{thm:sheaf}
$\mathrm{banachPresheaf}\;V$ is a sheaf for the coherent topology on $\CompHaus$.
\end{theorem}

The proof transfers the sheaf condition from the Type-valued \texttt{yonedaPresheaf} through the forgetful functor. The main technical obstacle was universe arithmetic, bridged using \texttt{hasLimitsOfSizeShrink} and \texttt{preservesLimitsOfSize\_of\_univLE}.

\begin{theorem}[Sorry-free]
There exists a functor $\SemiNormedGrp \to \CondAb$.
\end{theorem}

\subsection{Session 5: Faithful but Not Full}

The fifth session established the central structural property of the embedding, and did so at two levels.

\begin{theorem}[Sorry-free]
The sheaf-level embedding $\SemiNormedGrp \to \CondAb$ is faithful.
\end{theorem}

The faithfulness proof is machine-checked: constant functions at the one-point space separate the points of any seminormed group, so distinct bounded maps induce distinct condensed morphisms.

To machine-check non-fullness as well, we work one categorical layer down, with the presheaf $S \mapsto C(S, A)$ valued in abelian groups, assembled into a functor
\[
E : \SemiNormedGrp \longrightarrow \mathrm{PSh}(\mathrm{CompHaus}, \mathrm{Ab}).
\]
This is the presheaf underlying the condensed embedding; passing to it costs only the sheaf condition, which is irrelevant to the fullness obstruction (a statement purely about morphisms).

\begin{theorem}[Sorry-free]
\label{thm:fullness}
The functor $E$ is faithful but \textbf{not} full. Both statements are machine-checked for the same functor.
\end{theorem}

\begin{proof}
Faithfulness is the injectivity of postcomposition on bounded maps, verified by evaluation at the one-point space. For non-fullness, let $V = \mathbb{N} \to_0 \mathbb{Z}$ (finitely supported integer sequences) with the supremum norm. Because every nonzero integer entry has absolute value at least $1$, the norm satisfies $\|a\| \ge 1$ for $a \neq 0$, so $V$ carries the \emph{discrete} topology; consequently every additive map out of $V$ is continuous. The summation map $f(a) = \sum_n a_n$ is therefore a continuous additive homomorphism, but it is not bounded: the indicators $w_N$ of $\{0, \dots, N\}$ satisfy $\|w_N\| = 1$ while $f(w_N) = N+1$. Hence the natural transformation $E$ assigns to $f$ has no preimage among bounded maps, and $E$ is not full.
\end{proof}

The discreteness of $V$ is the engine: the presheaf realization sees continuous additive maps while $\SemiNormedGrp$ admits only bounded ones. The counterexample theorem itself contains no placeholder.

\textbf{From presheaves to sheaves.} The sheaf condition was already proved in \texttt{BanachEmbedding.lean}. The file \texttt{SheafFullnessCounterexample.lean} universe-lifts the summation example and packages it as a genuine morphism in $\CondAb$, yielding a sorry-free theorem that \texttt{semiNormedGrpToCondensedAb} is not full.

\textbf{Root cause.} Continuous additive maps between seminormed abelian groups need not be bounded. The standard proof that ``continuous linear implies bounded'' requires scalar multiplication by a dense subfield. Over real normed spaces the immediate boundedness obstruction disappears, but categorical fullness still needs proof. Over complex or general fields, $\CondAb$ forgets scalar linearity, so a condensed-module target is required.

The session also supplied a sorry-free finite-products instance for $\SemiNormedGrp$ in the pinned Mathlib version.

\subsection{Session 6: Frobenius Presentation}

The sixth session (\texttt{Cob2.lean}) formalized commutative Frobenius algebra data and an ordinary presentation inspired by the classical 2-dimensional classification~\cite{Kock2003}. It does not prove that classification.

\begin{definition}[Sorry-free]
A \emph{Frobenius object} in a monoidal category extends Mathlib's monoid and comonoid objects with the Frobenius relation:
\[
(\mu \otimes \id) \circ (\id \otimes \delta) = \delta \circ \mu.
\]
A \emph{commutative Frobenius object} additionally requires $\mu \circ \beta = \mu$.
\end{definition}

\begin{theorem}[Sorry-free]
In any braided monoidal category, the unit object carries a canonical commutative Frobenius algebra structure.
\end{theorem}

An ordinary Frobenius presentation category is defined on the natural numbers, with morphisms generated by $\mu$ (pants), $\eta$ (cap), $\delta$ (copants), $\varepsilon$ (cup), plus composition, tensor, and swap, quotiented by the Frobenius relations.

%% ============================================================
\section{Results Summary}
\label{sec:results}
%% ============================================================

\subsection{Theorem Inventory}

\begin{table}[h]
\centering
\small
\begin{tabular}{@{}lll@{}}
\toprule
\textbf{Result} & \textbf{File} & \textbf{Status} \\
\midrule
$\CondAb$ is abelian & LiquidTQFT.lean & Sorry-free \\
Mono + epi = iso & LiquidTQFT.lean & Sorry-free \\
$\mathrm{MonoidalCategory}\;\CondAb$ & LiquidTQFT.lean & Sorry-free \\
$\mathrm{BraidedCategory}\;\CondAb$ & LiquidTQFT.lean & Sorry-free \\
$\mathrm{SymmetricCategory}\;\CondAb$ & LiquidTQFT.lean & Sorry-free \\
\texttt{AbstractTQFT.transfer} & LiquidTQFT.lean & Sorry-free \\
$W.\mathrm{IsMonoidal}$ for coherent topology & MonoidalViaLoc.lean & Sorry-free \\
\texttt{presheafToSheaf} is monoidal & MonoidalViaLoc.lean & Sorry-free \\
\texttt{banachPresheaf} $V$ is a sheaf & BanachEmbedding.lean & Sorry-free \\
$\SemiNormedGrp \to \CondAb$ functor & BanachEmbedding.lean & Sorry-free \\
Embedding is faithful (sheaf level) & BanachEmbedding.lean & Sorry-free \\
Embedding $E$ faithful, not full (presheaf) & FullnessCounterexample.lean & Sorry-free \\
Sheaf-level embedding not full & SheafFullnessCounterexample.lean & Sorry-free \\
$\SemiNormedGrp.\mathrm{hasFiniteProducts}$ & BanachEmbedding.lean & Sorry-free \\
\texttt{FrobeniusObj}, \texttt{CommFrobeniusObj} & Cob2.lean & Sorry-free \\
Trivial Frobenius on $\mathbf{1}_C$ & Cob2.lean & Sorry-free \\
Frobenius presentation category (ordinary only) & Cob2.lean & Sorry-free \\
\bottomrule
\end{tabular}
\caption{Theorem inventory for the machine-checked development. Non-fullness is verified both for the auxiliary presheaf functor and for the genuine sheaf-level functor into $\CondAb$.}
\label{tab:results}
\end{table}

\subsection{Remaining Sorries}

Five sorry placeholders remain across three incomplete results:
\begin{enumerate}
\item \texttt{PreservesEqualizers} for the embedding (hard; 4--5 layers of limit transfer).
\item \texttt{PreservesFiniteProducts} for the embedding (hard; \texttt{ContinuousMap.piEquiv} path).
\item \texttt{CommFrobeniusData.toCob2Functor} (tedious; three placeholders for the functor's action on morphisms and its two functoriality proofs, all blocked on case analysis over the quotient).
\end{enumerate}

\subsection{Key Findings}

\textbf{Finding 1.} The monoidal structure on $\CondAb$ was already in Mathlib. The iterative process (axiom $\to$ manual construction $\to$ localization discovery) demonstrates the difficulty of navigating large formal libraries.

\textbf{Finding 2.} No packaged bridge between the relevant normed-group and condensed components was found in Mathlib v4.28.0. This is a version-scoped observation, not a global priority claim.

\textbf{Finding 3.} Fullness fails for both the auxiliary presheaf functor and the sheaf-level embedding, machine-checked by the summation counterexample. The condensed setting sees strictly more morphisms than the normed category, because the sup-norm topology on integer sequences is discrete and admits unbounded continuous additive maps.

\textbf{Finding 4.} The project supplies $\SemiNormedGrp.\mathrm{hasFiniteProducts}$ for the pinned Mathlib version.

\textbf{Finding 5.} The commutative Frobenius algebra data used in the classical classification~\cite{Kock2003} formalizes cleanly from existing \texttt{Mon\_}/\texttt{Comon\_} infrastructure.

%% ============================================================
\section{Roadmap for Further Development}
\label{sec:roadmap}
%% ============================================================

\subsection{Near-term}
Preserve finite limits for the embedding; complete the ordinary presentation functor; then formalize monoidal coherence and its universal property.

\subsection{Medium-term}
Define scalar-sensitive normed-space and condensed-module categories and state a correctly hypothesized fullness problem; construct a concrete Frobenius object in $\CondAb$; formulate a precise analytic exactness theorem with explicit hypotheses.

\subsection{Long-term}
Formalize a suitable projective tensor product; identify and port a precise LTE-related theorem with its original hypotheses; formalize genuine bordism categories.

\subsection{Frontier}
Concrete Chern--Simons constructions valued in $\CondAb$; Costello--Gwilliam~\cite{CostelloGwilliam2017} factorization algebras in the condensed setting.

%% ============================================================
\section{Methodology}
\label{sec:methodology}
%% ============================================================

The work was conducted in six sessions of AI-assisted formal verification. Each session consisted of: (1)~a carefully scoped prompt from Claude to Aristotle, based on analysis of previous results; (2)~Aristotle producing a Lean~4 file with Mathlib API audit, constructions, and proofs; (3)~Claude analyzing the output and designing the next prompt.

The prompts evolved from exploratory to targeted. The most productive prompts front-loaded reconnaissance (searching Mathlib for existing infrastructure) before attempting construction. The Session~3 breakthrough directly resulted from this ``search before build'' approach.

\subsection{Tools}
\begin{itemize}
\item \textbf{Lean~4} with \textbf{Mathlib v4.28.0}~\cite{Mathlib} for formal verification.
\item \textbf{Harmonic Aristotle}~\cite{Aristotle2025} for proof search and Mathlib API discovery.
\item \textbf{Claude (Anthropic)} for mathematical synthesis, prompt engineering, and documentation.
\item The project compiles with \texttt{lake build}; build warnings identify the five remaining placeholders.
\end{itemize}

%% ============================================================
\section{Conclusion}
\label{sec:conclusion}
%% ============================================================

The project contains six active Lean files with five placeholders and no custom axiom declarations. The central verified results are the Mathlib-supplied symmetric monoidal structure on $\CondAb$, braided monoidal composition scaffolding, the sheaf realization of seminormed groups, sorry-free faithfulness and non-fullness theorems for that realization, finite products in $\SemiNormedGrp$, and commutative Frobenius algebra data.

A notable result is the machine-checked sheaf-level non-fullness theorem. Universe-lifted finitely supported integer sequences have a discrete sup-norm topology, so summation is continuous but unbounded; it therefore defines a condensed morphism that cannot arise from a bounded morphism in $\SemiNormedGrp$.

The remaining barriers are specific: choosing and formalizing an analytic source category and tensor product, stating any LTE-related theorem with precise hypotheses, and constructing a genuine symmetric monoidal bordism presentation. Applications to gauge theory or factorization algebras remain research directions rather than established consequences of this codebase.

\subsection*{Code Availability}

The current Lean~4 project is available at \url{https://github.com/seanm27lol/Liquid-tft-lean-CM4CM-} and builds against Mathlib v4.28.0.

\subsection*{Acknowledgments}

Formal verification assisted by Aristotle~\cite{Aristotle2025} (Harmonic). Mathematical synthesis and prompt engineering by Claude (Anthropic). The symmetric monoidal structure on $\CondAb$ assembled in this project relies entirely on Mathlib infrastructure built by Jo\"el Riou and Dagur Asgeirsson. We thank Ya\"el Dillies and Jo\"el Riou for feedback on the Lean Zulip.

%% ============================================================
\bibliographystyle{amsalpha}
\bibliography{references}

\end{document}
